\subsection{FRM Exam 2009: which multi-day VaR is inconsistent?}

\textbf{Enunciado}

\emph{Assume that portfolio daily returns are independent and identically normally distributed. Which one of the following VARs is inconsistent with the others?}
\begin{enumerate}[label=\alph*.]
\item \hl{VAR(10-day) = USD 316M}
\item VAR(15-day) = USD 465M
\item VAR(20-day) = USD 537M
\item VAR(25-day) = USD 600M
\end{enumerate}

\subsubsection*{Solución}


\textbf{Propiedad de escala bajo i.i.d. normal}. Como se menciona antes, la varianza escala linealmente con el tiempo y la desviación estándar como \(\sqrt{h}\). Para un mismo \(\alpha\): \(\VaR(h)\propto \sqrt{h}\). Entonces \(\VaR(h)/\sqrt{h}\) debe ser constante.

\textbf{Normalizar cada número}\\
\[
316/\sqrt{10}\approx 100,\qquad
465/\sqrt{15}\approx 120,\qquad
537/\sqrt{20}\approx 120,\qquad
600/\sqrt{25}=120.
\]
Tres valores implican una VaR diaria consistente \(\approx 120\); el de 10 días implica \(\approx 100\). Por ello el inconsistente es el de 10 días.
